\documentclass[aspectratio=169,8pt]{beamer}
\usetheme[
%%% options passed to the outer theme
%    progressstyle=movCircCnt,   %either fixedCircCnt, movCircCnt, or corner
%    rotationcw,          % change the rotation direction from counter-clockwise to clockwise
%    shownavsym          % show the navigation symbols
  ]{AAUsimple}
  
% If you want to change the colors of the various elements in the theme, edit and uncomment the following lines
% Change the bar and sidebar colors:
%\setbeamercolor{AAUsimple}{fg=red!20,bg=red}
%\setbeamercolor{sidebar}{bg=red!20}
% Change the color of the structural elements:
%\setbeamercolor{structure}{fg=red}
% Change the frame title text color:
%\setbeamercolor{frametitle}{fg=blue}
% Change the normal text color background:
%\setbeamercolor{normal text}{fg=black,bg=gray!10}
% ... and you can of course change a lot more - see the beamer user manual.

\usepackage{fontspec}
\usepackage{unicode-math}
\setmathfont[version=bold]{XITS Math Bold}
\setmathfont{XITS Math}\setmainfont{Arial}
\setsansfont{Arial}
% \setsansfont{SourceSansPro-Regular.otf}

\usepackage{xeCJK}
\setCJKsansfont{FandolHei}
\setCJKmainfont{FandolSong}

% \usepackage[utf8]{inputenc}
\usepackage[english]{babel}
% \usepackage[T1]{fontenc}
% Or whatever. Note that the encoding and the font should match. If T1
% does not look nice, try deleting the line with the fontenc.
% \usepackage{helvet} 
\usepackage{indentfirst}
\usepackage{tikz}
\usetikzlibrary{backgrounds,mindmap}
\usepackage{xcolor}
\usetikzlibrary{calc,positioning,intersections}
\usepackage{pgfplots}
\pgfplotsset{compat=1.18}
\usepackage{listings}
\usepackage{caption}
% \usepackage{enumitem} % 与 beamer 冲突会导致 labelenumi 递归报错，先不用
% \setlist[itemize]{itemsep=6ex,parsep=4ex} % item 之间与同 item 内段落/换行的间距

% for algorithm
% ---------------------------------------
\usepackage{latexsym}
\usepackage{amsmath,amssymb,bm}
\usepackage{color,xcolor,tikz}
\usepackage{graphicx}
\usepackage{algorithm}
\usepackage{algorithmic}
\usepackage{amsthm}
\usepackage{float}
\usepackage{setspace}
\usepackage{hyperref}
\hypersetup{hidelinks}
% \usepackage[linesnumbered,vlined,boxed,commentsnumbered]{algorithm2e}
\usepackage{pdflscape}
% ---------------------------------------

% \definecolor{blockTColor}{RGB}{18, 67, 180}
% \definecolor{blockTColor}{RGB}{239, 131, 84} % coral orange
% \definecolor{blockTColor}{RGB}{7, 71, 143} % light blue
% \setbeamercolor{block title}{bg=blockTColor,fg=white}
% \definecolor{blockBColor}{RGB}{248, 247, 255}
% \setbeamercolor{block body}{bg=blockBColor,fg=black}
% \setbeamercolor{block title}{bg=cyan!90,fg=white}
% \setbeamercolor{block body}{bg=blue!10,fg=black}
% \useoutertheme{infolines}
% colored hyperlinks
\newcommand{\chref}[2]{%
  \href{#1}{{\usebeamercolor[bg]{AAUsimple}#2}}%
}

\title{An Improved Syndrome Bit Flipping Decoder with Narrowed Flipping Set for LDPC Codes}

% \subtitle{ICCC 2023}  % could also be a conference name
\date{December 9, 2023\\}
\author{
  Wenlin Zhao, Jie Zhang, Li Li and Haitong Jiang
}

% - Give the names in the same order as they appear in the paper.
% - Use the \inst{?} command only if the authors have different
%   affiliation. See the beamer manual for an example

\institute[
%  {\includegraphics[scale=0.2]{aau_segl}}\\ %insert a company, department or university logo
  School of Information Science and Technology\\
  Southwest Jiaotong University\\
  Cheng du, China
] % optional - is placed in the bottom of the sidebar on every slide
{% is placed on the bottom of the title page
  School of Information Science and Technology\\
  Southwest Jiaotong University\\
  Cheng du, China
  
  %there must be an empty line above this line - otherwise some unwanted space is added between the university and the country (I do not know why;( )
}

% specify a logo on the titlepage (you can specify additional logos an include them in 
% institute command below
\pgfdeclareimage[height=1.8cm]{titlepagelogo}{AAUgraphics/swjtu_circle_light_blue} % placed on the title page
\pgfdeclareimage[height=1.8cm]{titlepagelogo2}{AAUgraphics/ICCC_logo} % placed on the title page
\titlegraphic{% is placed on the bottom of the title page
  \pgfuseimage{titlepagelogo}
  \hspace{0.75cm}\pgfuseimage{titlepagelogo2}
}

\begin{document}
% the titlepage
{\aauwavesbg%
\begin{frame}[plain,noframenumbering] % the plain option removes the header from the title page
  % \fontsize{12}{12}\selectfont
  \titlepage
\end{frame}}
%%%%%%%%%%%%%%%%

% TOC
\begin{frame}{Tabel of Contents}{}
\tableofcontents
\end{frame}
%%%%%%%%%%%%%%%%

\section{SiDTBF: Merging Soft Info with DTBF}
\begin{frame}{Tabel of Contents}{}
	\frametitle{}
	\tableofcontents[currentsection]
\end{frame}

\section{DTBF Algorithm Overview}

\begin{frame}{Dynamic Threshold Bit Flipping (DTBF)}

  \begin{block}{Flipping Function ($E_n$)}
    The reliability $E_n$ of bit $n$ is calculated using the syndrome weight and a weighting factor:
    \begin{equation*}
      E_n = \alpha(z_n \oplus y_n) + \sum_{m \in A(n)}(2s_m - 1)
    \end{equation*}
    \vspace{-0.3cm}
    \begin{itemize}
      \item \textbf{Core Mechanism}: It introduces a dynamic threshold ($Th$) that adjusts during iterations to prevent the decoder from getting stuck in local loops when no bits satisfy the flipping condition.
      \item \textbf{Rule}: If $E_n \ge Th$, bit $n$ is flipped.
    \end{itemize}
  \end{block}

  \begin{block}{Dynamic Threshold Adjustment}
    \begin{itemize}
      \item \textbf{Initialization}: Start with maximum threshold $Th = d_v + \alpha$.
      \item \textbf{Adjustment}: Count bits ($t$) where $E_n \ge Th$.
      \item If $t=0$ (no bits flipped), reduce $Th$ by a step size $\Delta$ to escape the local deadlock.
      \item \textbf{Limitation}: Cannot effectively correct errors when RBER exceeds the algorithm's capability (e.g., $7 \times 10^{-3}$) as it treats all bits with equal initial probability.
    \end{itemize}
  \end{block}
\end{frame}

\begin{frame}{Soft Info from Additional Reads}
  \vspace{-0.1cm}
  \begin{columns}
    \begin{column}{0.55\textwidth}
      \begin{block}{Bound-based labeling}
        \begin{itemize}
          \item Insert two extra references in each overlap region; define \textbf{bound} $=y_1/y_2$ between adjacent PDFs.
          \item Cells between the added references are marked \textcolor{blue}{weak} ($b_i=0$); others are \textcolor{blue}{strong} ($b_i=1$).
          \item Bound must balance coverage of real errors vs. accidental inclusion of correct bits.
        \end{itemize}
      \end{block}
    \end{column}
    \begin{column}{0.4\textwidth}
      \centering
      \includegraphics[width=0.95\linewidth]{AAUgraphics/SiDTBF_multi_read.png}
      \\
      \tiny Threshold-voltage overlap; shaded area becomes weak bits.
    \end{column}
  \end{columns}
\end{frame}

\begin{frame}{Workflow and First-Iteration Logic}
  \vspace{-0.15cm}
  \begin{columns}
    \begin{column}{0.53\textwidth}
      \begin{block}{Algorithm sketch}
        \begin{enumerate}
          \setlength{\parsep}{4pt}\setlength{\itemsep}{4pt}
          \item Input hard bits $z$ plus weak/strong masks $B_{lsb/msb}$ (weak when sample lies between added refs).
          \item Iteration 1: compute sydrome weight $s$; compute $E_n=\alpha(z_n\oplus v_n)+\sum_{m\in\mathcal{A}(n)}(2s_m-1)$.
          \item Weak bits: if $E_n > E_n^{*}$ flip aggressively; strong bits still compare to $Th$ (DTBF rule).
          \item Iterations $>1$: drop to DTBF to prevent false-flip amplification.
          \item Stop when $s=\vec{0}$ or iter hits cap; output $y$.
        \end{enumerate}
      \end{block}
      \begin{block}{Rationale}
        \begin{itemize}
          \item Overlap-region errors dominate; lowering $E_n^{*}$ only in iter-1 cleans them quickly.
          \item Returning to DTBF keeps complexity intact and avoids later mis-flips.
        \end{itemize}
      \end{block}
    \end{column}
    \begin{column}{0.42\textwidth}
      \centering
      \includegraphics[width=0.9\linewidth]{AAUgraphics/SiDTBF_workflow.png}
      \\
      \tiny SiDTBF workflow
    \end{column}
  \end{columns}
\end{frame}

\begin{frame}{Tuning Bound and $E_{n}^{*}$}
  \vspace{-0.2cm}
  \begin{columns}
    \begin{column}{0.5\textwidth}
      \begin{block}{How thresholds are chosen}
        \begin{itemize}
          \item $E_n^{*}\!\in\!\{-d_v-\alpha+k\Delta\}$ with $\alpha\!=\!2,\Delta\!=\!1$; sweep across retention \& P/E conditions.
          \item Keep only best $E_n^{*}$ values for key retention points to control complexity.
          \item Bound is swept via chip scans and Gaussian fits; choose the value that maximizes corrected weak bits while limiting false flips.
        \end{itemize}
      \end{block}
    \end{column}
    \begin{column}{0.45\textwidth}
      \centering
      \includegraphics[width=0.95\linewidth]{AAUgraphics/SiDTBF_PE.png}
      \\
      \tiny Flipped-error ratio vs. $E_n^{*}$ across P/E cycles.
    \end{column}
  \end{columns}
\end{frame}

\begin{frame}{Performance Highlights}
  \vspace{-0.25cm}
  \begin{block}{Setup}
    \begin{itemize}
      \item 3D TLC NAND, $N=32768$, rate $0.89$, $d_v=5$, first-iter $E_n^{*}=-d_v+4$.
      \item Sweep bound $\in\{5,10,15,20,50\}$; bound $=20$ emerges best.
    \end{itemize}
  \end{block}
  \begin{columns}
    \begin{column}{0.52\textwidth}
      \begin{block}{Observed gains}
        \begin{itemize}
          \item Correctable RBER grows from $7\!\times\!10^{-3}$ (DTBF) to $9\!\times\!10^{-3}$.
          \item Bound $=20$ slashes FER from $6\!\times\!10^{-1}$ to $4\!\times\!10^{-5}$ at $6.8\!\times\!10^{-3}$.
          \item Average decoding latency drops $\approx3\times$ thanks to cleaner first iteration.
        \end{itemize}
      \end{block}
    \end{column}
    \begin{column}{0.43\textwidth}
      \centering
      \includegraphics[width=0.95\linewidth]{SiDTBF_Paper/images/9c0d4aef33a967a6b82c36b929b3a91925e97bd851809ca722ada35988b4a5ae.jpg}
      \\
      \tiny FER comparison.
    \end{column}
  \end{columns}
\end{frame}


\section{Context and Objectives}
\begin{frame}{Tabel of Contents}{}
	\frametitle{}
	\tableofcontents[currentsection]
\end{frame}

\begin{frame}{Context and Objectives}{}
\vspace{-0.3cm}

\begin{block}{Low Density Parity Check (LDPC) Decoder}
  \centering
  \includegraphics[width=0.75\textwidth]{/mnt/f/MyProject/LaTexProject/MaxPPT/SiDTBF/AAUgraphics/LDPC_Model.pdf}
\end{block}
\begin{block}{2 types of LDPC decoding algorithms:}
    \begin{itemize}
      \item \textbf{Soft information algorithms}: Sum-Product, Min-Sum...,\alert {powerful error correction} capacity but high complexity.
      \item \textbf{Hard-decision algorithms}: Bit Flipping (BF), Gradient Descent Bit Flipping (GDBF), etc. \alert{low complexity}, usually weak in error correction.
      \item \textbf{Towards a new type of decoder}: \alert{low complexity + randomness} \(\Leftrightarrow \) \alert{powerful error correction}.
    \end{itemize}
\end{block}
\end{frame}
%%%%%%%%%%%%%%%%

\begin{frame}{Trend and Experience}
  \begin{minipage}[c][.8\textheight][c]{.4\textwidth}
    \hspace{0.5cm}
    \includegraphics[width=0.6\textwidth]{/mnt/f/MyProject/LaTexProject/MaxPPT/SiDTBF/AAUgraphics/BF_Evolution_Timeline.pdf}
    % \captionof{figure}{BF Evolution}
  \end{minipage}%
  \hspace{-0.75cm} 
  \vspace{0.2cm}
  \begin{minipage}[c][.7\textheight][c]{.55\textwidth}

    \begin{block}{Trends of LDPC Bit Flipping decoders}
      \begin{enumerate}
        \itemsep1em
        \item \textbf{Injected Noise in the decoder can help to against the channel errors}: 
          \begin{itemize}
            \itemsep0.5em
            \item Probabilistic Gradient Descent Bit Flipping \alert{(PGDBF)} for BSC
            \item Noisy Gradient Descent Bit Flipping \alert{(NGDBF)} for AWGN
          \end{itemize}
        \item \textbf{Optimizing mechanism or strategy applyed on GDBF}:
          \begin{itemize}
            \itemsep0.5em
            \item Noisy Gradient Descent Bit-Flipping Decoder Based on Adjustment Factor \alert{(A-NGDBF)}
            \item Tabu-list Random-penalty Gradient Descent Bit Flipping \alert{(TRGDBF)}
            \item Gradient Descent Bit Flipping with Momentum \alert{(GDBF with Momentum)}
          \end{itemize}
        \item \textbf{Low complexity, as post-processing decoder}:
          \begin{itemize}
            \item Syndrome Bit Flipping \alert{(SBF)}
          \end{itemize}
      \end{enumerate}
      \par
  \vspace{\baselineskip}

   \end{block}
    % \begin{block}{Block 2}
    %   This is the second block on the right.
    % \end{block}
  \end{minipage}
\end{frame}


\begin{frame}{Proposed Algorithm}
  \vspace{-0.1cm}
  \centering
  \begin{minipage}{0.7\textwidth}
    % \begin{block}{}
      
      % \begin{tikzpicture}[remember picture, overlay]
      %   \node at (current page.west) [anchor=west, xshift=1cm] (leftbrace) {
      %     \scalebox{5}{$\left\lbrace\right.$}
      %   };
      % \end{tikzpicture}

        \textbf{Key Concepts:}
        \begin{itemize}
          \item Improved Syndrome Bit Flipping (SBF) Algorithm.
          \item Incorporation of a \textcolor{blue}{Narrowed Flipping Set (NFS)}.
        \end{itemize}

        \vspace{\baselineskip}

        \textbf{Algorithm Highlights:}
        \begin{itemize}
          \item Combines channel information and random selection for NFS.
          \item Mitigates error propagation during bit-flipping.
          \item Aims for improved decoding convergence.
        \end{itemize}

        \vspace{\baselineskip}

        \textbf{Expected Advancements:}
        \begin{itemize}
          \item Significant Bit Error Rate (BER) performance improvement.
          \item Estimated \alert{1.2 dB} BER gain over original SBF.
          \item Balances enhanced BER performance with \textbf{modest complexity increase}.
        \end{itemize}

        \vspace{\baselineskip}
  
        \textbf{Applicability:}
        \begin{itemize}
          \item Effective across various LDPC codes.
          % \item Robust under varying SNR conditions.
        \end{itemize}
  
    % \end{block}
  \end{minipage}
\end{frame}

\section{Introduction of Related Algorithms}

\begin{frame}{Tabel of Contents}{}
	\frametitle{}
	\tableofcontents[currentsection]
\end{frame}

\begin{frame}{GDBF Algorithm}
  \vspace{-0.35cm}
  \begin{block}{Concepts of Gradient Descent Bit Flipping (GDBF) Algorithm}
    \begin{itemize}
      \item Iterative propagation of binary information between 2 groups of processing units
      \begin{enumerate}
        \item \textbf{Check Nodes Units} (CNU): compute parity check equations
        \begin{equation*}
          c_m^{(l)}=\bigoplus_{v_n \in \mathcal{N}(c_m)} v_n^{(l)}
        \end{equation*}
        \item \textbf{Variable Nodes Units} (VNU): the VN value is flipped if the number of violated CN neighbors is too large.
      \end{enumerate}
      \item Requires the computation of an Energy/Inversion Function in order to  \textcolor{blue}{select the bit-flips}.
    \end{itemize}
  \end{block}

  \begin{block}{Energy function for the AWGN channel}
    \setlength{\belowdisplayskip}{0pt}
    \setlength{\belowdisplayshortskip}{0pt}
    \centering \alert{\textbf{low energy}} \(\Leftrightarrow\) \alert{\textbf{low reliability}}
    \begin{equation*}
      E_{k}^{(l)} = x_k^{(l)} y_k + \sum (1-2c_m^{(l)})
    \end{equation*}
    \begin{itemize}
      \item \(x_k^{(l)} y_k\):\textbf{ maximum likelihood (ML) part}, correct codeword yields maximum value
      \item \(\sum (1-2c_m^{(l)})\): \textbf{penalty part}, the sum of the bipolar syndromes of \(\mathbf{x}\)
    \end{itemize}
    \end{block}
\end{frame}


% \subseciton{Some related BFs}
\begin{frame}{Variant of GDBF}
  \vspace{-0.3cm}
  \begin{block}{Noisy Gradient Descent Bit Flipping and Syndrome Bit Flipping}
    \begin{itemize}
      \item \alert{NGDBF}: Noisy version of GDBF, noise \(q_k \sim N(0,\eta^2 N_0/2)\) are injected to help GDBF get out of local maximum.
      \item \alert{SBF}: The channel information is \alert{no longer} taken into consideration.
      \begin{equation*}
        E_{k}^{(l)} = \sum c_m^{(l)}
      \end{equation*}
    \end{itemize}

  \end{block}

  \vspace{-0.2cm}
  \begin{block}{Unified presentation of GDBF, NGDBF and SBF decoders}
    \textbf{[Step 1] Compute CNs values}: \(c_m^{(l)},\forall m=1,\ldots,M\) \\
    \textbf{[Step 2] Compute Energy function of VNs}: \(E_{k}^{(l)},\forall n=1,\ldots,N\)\\
    \textbf{[Step 3] Determine the flipping set}
    \begin{gather*}
      \mathcal{S}_{\text{GDBF}}^{(l)} = \{n\in\{1,N\};E_k^{(l)}\le \theta \} \\ 
      \mathcal{S}_{\text{NGDBF}}^{(l)} = \{n\in\{1,N\};E_k^{(l)}\le \theta + q_k^{(l)} \}\\ 
      \mathcal{S}_{\text{SBF}}^{(l)} = \{n\in\{1,N\};E_k^{(l)} > \theta_k(\ell) \} 
    \end{gather*}
    \textbf{[Step 4] Bit flipping}
    \begin{equation*}
      \forall n \in \mathcal{S}^{(l)}:  x_n = 1-x_n
    \end{equation*}
  \end{block}
\end{frame}

\begin{frame}{Flipping Sets}
  \vspace{-0.3cm}
  \begin{block}{Flipping sets contain indices of bits to be flipped}
    \vspace{-\baselineskip} % 可用于block中整体上移
    \begin{gather*}
      \mathcal{S}_{\text{GDBF}}^{(l)} = \{n\in\{1,N\};E_k^{(l)}\le \theta \} \\ 
      \mathcal{S}_{\text{NGDBF}}^{(l)} = \{n\in\{1,N\};E_k^{(l)}\le \theta + q_k^{(l)} \}\\ 
      \mathcal{S}_{\text{SBF}}^{(l)} = \{n\in\{1,N\};E_k^{(l)} > \theta_k(\ell) \} 
    \end{gather*}
  \end{block}
  \vspace{-0.1cm}
  \begin{block}{Graphical representation of GDBF and NGDBF flipping sets}
    \centering
    \includegraphics[width=0.7\textwidth]{/mnt/f/MyProject/LaTexProject/MaxPPT/SiDTBF/AAUgraphics/Flipping_Set_Graphic.pdf}
  \end{block}
\end{frame}
%%%%%%%%%%%%%%%%

\begin{frame}{SBF Algorithm}
  \vspace{-0.3cm}
  \begin{block}{Why \(E_{k}^{(l)} = \sum c_m^{(l)}\) works?}
    \vspace{\baselineskip}
    \centering \alert{Lower accuracy} \(\rightarrow\) \alert{Concise Architecture}
    \begin{itemize} 
      \item Old idea\footnote[frame]{Raul Benet, Adriaan De Lind Van Wijngaarden, Ralf Dohmen, Thomas Richardson, Rudiger
      Urbanke, “Iterative decoding of low-density parity-check (LDPC) codes”, Patent EP1643653A1, 5
      april 2006.}: Use the number of unsatisfied checks connected to a VN to
      represent the reliability of VN.\\
      \item \textbf{Syndrome can be viewed as a compressed representation of the error pattern}.

    \end{itemize}
    
    % \footnote[frame]{Second footnote.}. 可在页面底部生成footnote
    \includegraphics[width=0.8\textwidth]{/mnt/f/MyProject/LaTexProject/MaxPPT/SiDTBF/AAUgraphics/SBF_E_Graphic.eps}
  \end{block}
\end{frame}


\begin{frame}{Example of a SBF decoding processing}
  \vspace{-0.2cm}
  \centering
  \includegraphics[width=0.8\textwidth]{/mnt/f/MyProject/LaTexProject/MaxPPT/SiDTBF/AAUgraphics/SBF_E_over4.pdf}
  \begin{columns}
    \hspace{0.2cm}
    \begin{column}{0.45\textwidth}
      Flip bits with \(E > 4\) \(\Rightarrow \) 19 bits flipped :\\
      \hspace*{11em} 14 that was in error \\ % Adjust the space as needed
      \hspace*{11em} \alert{5 that was correct}\\
      Only 20 – 14 + 5 = 11 errors after the first iteration.\\
      \includegraphics[width=0.7\textwidth]{/mnt/f/MyProject/LaTexProject/MaxPPT/SiDTBF/AAUgraphics/SBF_process.pdf}
    \end{column}
    \begin{column}{0.4\textwidth}
      \vspace{-0.3cm}
      \begin{block}{Threshold sequence \(\theta\)}
        A threshold sequence is << \alert{a decoding key} >>.
        \begin{itemize}
          \item \textbf{Non-unique}.
          \item Searched by \textbf{trapping set analysis}.
          \item \(0\le\theta\le\max(d_v)\).
        \end{itemize}
      \end{block}
    \end{column}
  \end{columns}
\end{frame}

\section{The Narrowed Flipping Set}

\begin{frame}{Tabel of Contents}{}
	\frametitle{}
	\tableofcontents[currentsection]
\end{frame}

\begin{frame}{The unnecessary flip}
  \begin{block}{Potential problem during bit flipping}
    All ‘0’ codeword transmission:
    \begin{itemize}
      \item Non-error bits are \textbf{unnecessarily flipped} in threshold-based bit flipping algorithms.
    \end{itemize}
    \centering
    \includegraphics[width=0.75\textwidth]{/mnt/f/MyProject/LaTexProject/MaxPPT/SiDTBF/AAUgraphics/unneccessary_flip_SBF.pdf}
  \end{block}
\end{frame}

\begin{frame}{The unnecessary flip}
  \begin{block}{Potential problem during bit flipping}
    All ‘0’ codeword transmission:
    \begin{itemize}
      \item Non-error bits are \textbf{unnecessarily flipped} in threshold-based bit flipping algorithms.
    \end{itemize}
    \centering
    \includegraphics[width=0.75\textwidth]{/mnt/f/MyProject/LaTexProject/MaxPPT/SiDTBF/AAUgraphics/unneccessary_flip_MNGDBF.pdf}
  \end{block}

\end{frame}

\begin{frame}{An underdeveloped idea from paper}

  \begin{columns}
    \hspace{0.5cm}
    \begin{column}{0.4\textwidth}

      \begin{block}{Heuristic in SBF paper}
        \setstretch{1.25}
        \begin{itemize}
          \item \textbf{Heuristic 3}\footnote[frame]{E. Boutillon, C. Winstead, and F. Ghaffari, “The Syndrome Bit Flipping Algorithm for LDPC Codes,” IEEE Communications Letters, pp. 1–1, 2023}. : \textcolor{blue}{SBF is only applied for bits with "unreliable" channel samples}, $|y|\le\gamma$ for some threshold $\gamma$. The choice of \(\gamma\) is determined empirically by simulation.
        \end{itemize}
      \end{block}


      \begin{block}{The limitaion of Heuristic 3}
        \begin{itemize}
          \item BER gain is small.
          \item Optimal $\gamma$ is varying with SNR.
        \end{itemize}
      \end{block}
    \end{column}
    \begin{column}{0.5\textwidth}

        \centering
        \includegraphics[width=0.9\textwidth]{/mnt/f/MyProject/LaTexProject/MaxPPT/SiDTBF/AAUgraphics/SFB+h3.jpg}

    \end{column}
  \end{columns}
\end{frame}

\begin{frame}{Narrow down the Flipping Set}
  \vspace{-0.4cm}
  \begin{block}{The experiment for narrowing down the flipping set}
    Construct a flipping set \(\mathcal{S}\) that includes the positions of all error bits \(\mathcal{S}_{\text{err}}\) and another \(n\) non-error bits.
    \begin{itemize}
      \item During SBF decoding, flipping bits are selected in \(\mathcal{S}\).
      \item The original SBF algorithm: a circumstance where \(\mathcal{S}\) contains all the bits in the codeword.
      % \item as \(n\) decreases, the flipping set \(\mathcal{S}_{\text{SBF}}\) becomes increasingly constrained by \(\mathcal{S}\).
    \end{itemize}
  \end{block}
  \vspace{0.1cm}
  \begin{columns}
    \begin{column}{0.5\textwidth}
      % \begin{block}{}
        \centering
        \includegraphics[width=0.75\textwidth]{/mnt/f/MyProject/LaTexProject/MaxPPT/SiDTBF/AAUgraphics/SBF_with_Fn.pdf}
      % \end{block}
    \end{column}
    \begin{column}{0.35\textwidth}
      \begin{block}{}
        \begin{itemize}
          \item \textcolor{blue}{\textbf{Conclusion}}: Narrowing down the flipping set makes SBF decoding more convergent.
        \end{itemize}
      \end{block}
    \end{column}
  \end{columns}
\end{frame}


\begin{frame}{Narrowed Flipping Set}
  \vspace{-0.1cm}
  \begin{block}{Construct the flipping set \(\mathcal{S}\)}
    \centering \alert{Channel Information + Random Selection}
    \begin{itemize}
      \item Bits with lowest magnitude of \(|⃗\vec{y}|\) are considered as unreliable bits. 
    \end{itemize}

    \vspace{0.3em}

    \begin{columns}%[T] % T aligns the tops of the columns
      \begin{column}{0.5\textwidth}
        \includegraphics[width=0.8\textwidth]{/mnt/f/MyProject/LaTexProject/MaxPPT/SiDTBF/AAUgraphics/bits_include_probability.pdf}
        % \tiny Monte Carlo simulation of the average ratio of the number of error bits contained within the \textbf{top-n unreliable bits} to the total number of error bits.
      \end{column}
      \begin{column}{0.4\textwidth}
        \begin{equation*}
          P_n = \frac{\text{error bits within }n \text{ smallest } |\vec{y}|}{\text{total error bits}} 
        \end{equation*}
        \begin{equation*}
          \mathcal{S}=\mathcal{S}_{CI}\cup \mathcal{S}_{\text{Random}}
        \end{equation*}
        
        \begin{itemize}
          \item \(\mathcal{S}_{\text{err}} \subseteq \mathcal{S}\) is essential for successful decoding.
          \item Flipping bits are selected in \(\mathcal{S}_{\text{SBF}} \cap \mathcal{S}\).
        \end{itemize}
      \end{column}
    \end{columns}
  \end{block}
\end{frame}

% \begin{landscape}
\begin{frame}{Algorithm}
  \centering
  \vspace{-0.4cm}
  \begin{minipage}{0.6\textwidth}

  \begin{algorithm}[H]
  \caption{SBF decoding with narrowed flipping set}
  \begin{algorithmic}  % enter the algorithmic environment
  {\small 
    \REQUIRE  $\vec{y} = (y_1,\dots,y_N)\in{\mathbb R}^N$ $\quad${(received word)} \\
                $\qquad \quad \Theta_k = \{\theta_k(\ell)\}_{\ell=1\ldots N_k}$ $\qquad${(threshold sequences)}

    % \STATE $m=0;\vec{x}=(1-\operatorname{sign}(\vec{y}))/2$;
    % \STATE $\vec{x}_0 = \vec{x};\vec{s}_0=\mathbf{H}\vec{x}_0 \; \operatorname{mod}2$;
    \medskip

    \STATE $\mathcal{S}_1 \gets  T_1 \; \text{smallest} \; |\vec{y}| \text{ bit positions}; \qquad{\rhd \textbf{\textcolor{blue}{Channel information part}}}$
    \WHILE {$m < \gamma \; \textbf{and} \; \vec{s} \neq \vec{0} $}

      \STATE $k=0;m=m+1$;
      \STATE $\mathcal{S}_2 \gets \textcolor{blue}{\text{random select }} T_2 \text{ bit positions from } \vec{y} \setminus  \mathcal{S}_1; \qquad{\rhd \textbf{\textcolor{blue}{Random part}}}$
      \STATE $\mathcal{S} \gets \mathcal{S}_1 \cup \mathcal{S}_2;$
      \IF {$m = 2$}
        \STATE 	$\Theta = \Theta(1); \qquad{\rhd \textbf{\textcolor{blue}{Reduce computational complexity}}}$
      \ENDIF

      % \WHILE {$ k<\operatorname{size}(\Theta) \; \textbf{and} \; \vec{s} \neq \vec{0} $} %$\quad${\textbf{SBF iteration loop}}
      %   \STATE $\ell = 0;k=k+1;\vec{}=\vec{s}_0;\vec{x}=\vec{x}_0$;

      %   \WHILE {$\ell<N_k \; \textbf{and} \; \vec{s} \neq \vec{0} $}
      \WHILE {\textcolor{blue}{SBF Loop}}
          % \STATE $\ell = \ell + 1$;
          \STATE $E = \mathbf{H}^T \vec{s}$; %$\quad${energy computation}
          \STATE $\mathcal{S}_{\text{SBF}} \gets \{n\in\{1\ldots N\}; E_n>\theta_k(\ell)\}$;
          \IF{$m=1$} 
            \STATE $\forall n \in \mathcal{S}_{\text{SBF}} : x_n = 1- x_n; \qquad\qquad \rhd \textbf{\textcolor{blue}{SBF iteration}}$
          \ELSE
            \STATE $\forall n \in \mathcal{S}_{\text{SBF}}\cap \mathcal{S} : x_n =  1- x_n; \qquad \rhd \textbf{\textcolor{blue}{NFS-SBF iteration}}$
          \ENDIF
          \STATE $\vec{s}=\mathbf{H}\vec{x} \; \operatorname{mod}2$;
        \ENDWHILE
      \ENDWHILE
    % \ENDWHILE
    % \smallskip
    % \IF {$\vec{s} \neq \vec{0}$}
    %   \STATE $\vec{x}=\vec{x}_0$;
    % \ENDIF
  }
  \end{algorithmic}
  \end{algorithm}

\end{minipage}
\end{frame}
% \end{landscape}

\section{Parameter Optimizing and Performance Evaluation}
\begin{frame}{Tabel of Contents}{}
	\frametitle{}
	\tableofcontents[currentsection]
\end{frame}

\begin{frame}{Parameter Optimizing}
  \vspace{-0.4cm}
  \begin{block}{}
    We denote $T_1,T_2$ as size of $\mathcal{S}_{CI},\mathcal{S}_{\text{Random}}$, there are some facts:
    \begin{itemize}
      \item The probability $p$ that of $\mathcal{S}_{\text{err}} \subseteq \mathcal{S},(\mathcal{S}=\mathcal{S}_{CI}\cup \mathcal{S}_{\text{Random}})$ increases with the growth of $T_1+T_2$.
      \item BER performance degrades with the growth of $T_1+T_2$.
      \item $T_1+T_2$ is expected as low as possible.
    \end{itemize}
  \end{block}

  \vspace{-0.2cm}
  \begin{columns}%[T] % T aligns the tops of the columns
    \begin{column}{0.4\textwidth}
      \begin{itemize}
        \setstretch{1.5}
        \item Rest of error bits after $(T_1,T_2)$ extraction:\\
              $N_r = N\cdot P_e\cdot(1-P_n)$
        \item $p = \frac{\binom{N-T_1-N_r}{T_2-N_r} }{\binom{N-T_1}{T_2} }.$
        \item Generalized binomial coefficient: $\binom{n}{k} = \frac{\Gamma(n+1)}{\Gamma(k+1)\Gamma(n-k+1)}.$
        \item $p =\frac{\Gamma(N-T_1-N_r+1)\Gamma(T_2+1)}{\Gamma(T_2-N_r+1)\Gamma(N-T_1+1)}.$ 
      \end{itemize}
    \end{column}

    \begin{column}{0.5\textwidth}
      
      \setbeamerfont{block title}{size=\normalsize}
      \begin{block}{Probability \(p\) for different $T1$ and $T2$.}
        \centering
        \includegraphics[width=0.75\textwidth]{/mnt/f/MyProject/LaTexProject/MaxPPT/SiDTBF/AAUgraphics/p_T1T2.pdf}
      \end{block}
    \end{column}
  \end{columns}
\end{frame}

\begin{frame}{Optimal Selection of $T_1$ and $T_2$}

  \vspace{-0.2cm}
  \centering
  \begin{minipage}{0.9\textwidth}
  \begin{block}{$T_1$: Channel Information Utilization}
    \begin{itemize}
      \item Purpose: Identify the most unreliable bit positions.
      \item Strategy: Select top $T_1$ bits with smallest absolute values.
    \end{itemize}
  \end{block}

  \begin{block}{$T_2$: Random Selection Complement}
    \begin{itemize}
      \item Purpose: Increase likelihood of encompassing all error bits.
      \item Strategy: Randomly select $T_2$ bits not in $T_1$.
    \end{itemize}
  \end{block}

  \begin{block}{SNR Considerations}
    \begin{itemize}
      \item High SNR: Smaller $T_1$ and $T_2$ may suffice due to concentrated errors.
      \item Low SNR: Larger $T_1$ and $T_2$ might be necessary for uniform error distribution.
    \end{itemize}
  \end{block}

  % \begin{block}{Experimental Tuning and Optimization}
  %   \begin{itemize}
  %     \item Method: Adjust based on simulation results and system performance.
  %     \item Goal: Optimal BER performance with balanced computational complexity.
  %   \end{itemize}
  % \end{block}
  \end{minipage}
\end{frame}


\begin{frame}{Performance Evaluation}
  \vspace{-0.3cm}
  \begin{block}{Performance of NFS-SBF}
    \begin{itemize}
      \item $(T_1,T_2)=(300,400)$ for 802.3an code, $(T_1,T_2)=(400,200)$ for PEG code.
      \item With growth of maximum $\gamma$ trials of $(T_1,T_2)$ extraction, NFS significantly improves the BER performance.
    \end{itemize}
    \centering
    \includegraphics[width=0.8\textwidth]{/mnt/f/MyProject/LaTexProject/MaxPPT/SiDTBF/AAUgraphics/performance_NFS_SBF.pdf}
  \end{block}
\end{frame}


\section{Conclusion and Heuristics}
\begin{frame}{Tabel of Contents}{}
  \hspace{5em}
	% \frametitle{}
	\tableofcontents[currentsection]
\end{frame}

\begin{frame}{Conclusion}
  \centering
  \begin{minipage}{0.7\textwidth}
  % \begin{block}{Narrowed Flipping Set (NFS)}
  %   \begin{itemize}
  %     \item Construct by channel information and random selection.
  %     \item Applicable for different LDPC codes.
  %     \item NFS successful avoids non-error bits from being unnecessarily flipped.
  %     \item The coding gains come at the cost of a modest increasing complexity.
  %   \end{itemize}
  % \end{block}


  \textbf{Main Conclusion:}
  \begin{itemize}
    \item Improves Bit Error Rate (BER) performance in LDPC codes.
    \item Integrates channel information and random selection for NFS.
    \item Mitigates error propagation in the bit-flipping process.
    \item Achieves better decoding convergence, especially at high SNR.
    \item Balances accuracy improvement with computational complexity.
  \end{itemize}
  \vspace{\baselineskip}
  \textbf{Key Takeaway:}
  \begin{itemize}
    \item Significant BER improvement over traditional SBF algorithms.
  \end{itemize}
  \vspace{\baselineskip}
  \textbf{Implications:}
  \begin{itemize}
    \item Validates NFS effectiveness in LDPC code decoding.
    \item Suggests potential for further error correction optimizations.
  \end{itemize}
\end{minipage}
\end{frame}

\begin{frame}{More Heuristics}
  \vspace{-0.1cm}
  \begin{block}{}
    \begin{itemize}
      \item Narrowed Flipping Set: \textbf{A cross-validation mechanism}.
      \item Can NFS be applied on other bit-flipping algorithms?-->\alert{YES}
      \begin{itemize}
        \item \textbf{Threshold-based bit flipping}: NGDBF, GDBF, MTBF
      \end{itemize}
    \end{itemize}

    \centering
    \includegraphics[width=0.4\textwidth]{/mnt/f/MyProject/LaTexProject/MaxPPT/SiDTBF/AAUgraphics/NFS_effect.jpg}
    
    \begin{itemize}
      \item The cost: Linearly increasing complexity.
    \end{itemize}
  \end{block}

\end{frame}


\begin{frame}{More Heuristics}
  \centering
  \begin{minipage}{0.9\textwidth}
    \begin{block}{Further Research}

        \textbf{The weakness of current NFS method:}
          \begin{itemize}
            \itemsep0.5em  % 在子列表中也可以设置
            \item When SNR is low, the distribution of error bit positions in the LDPC code becomes chaotic.
            \item The worst case come at the cost of a large number of iterations.
          \end{itemize}
        
        \vspace{\baselineskip}  

        \textbf{Further Research Directions:}
          \begin{itemize}
            \itemsep0.5em
              % \item \textbf{Refining NFS Construction}: Developing more effective methods to identify unreliable bit positions, potentially enhancing the algorithm's efficiency in lower SNR conditions.
              % \item \textbf{Iterative Improvement}: Exploring the use of information from previous iterations to optimize the NFS construction and potentially reduce the number of iterations required.
              % \item \textbf{Application to Other Algorithms}: Extending the NFS mechanism to other bit-flipping algorithms to evaluate its effectiveness and applicability across a broader range of decoding strategies and different LDPC codes.
              
              \item Developing more effective methods to identify unreliable bit positions.
              \item Leveraging information from previous iterations.
              \item Extending NFS mechanism to other algorithms.

          \end{itemize}

    \end{block}
  \end{minipage}
\end{frame}


% ------------------------------------------------------------
% Channel/HRE simulation appendix (HREModel project)
% ------------------------------------------------------------
\section{Channel Error Injection Models (ERR\_INJ / AWGN / TAWGN) and HRE}

\begin{frame}{ERR\_INJ: Fixed-Count Random Error Injection}
\vspace{-0.2cm}
\begin{columns}
  \begin{column}{0.56\textwidth}
    \begin{block}{What it models}
      \begin{itemize}
        \item Purpose: emulate a \textbf{fixed error weight} channel where the raw error count is controlled directly.
        \item Input: choose an integer $K$; randomly select $K$ bit positions; flip them deterministically.
        \item Output property: the frame-level error count is fixed: $FBC = K$ (by construction).
        \item This is useful for mapping \textbf{FER vs. error weight} without the randomness of AWGN tails.
      \end{itemize}
    \end{block}
    \begin{block}{How it differs from AWGN}
      \begin{itemize}
        \item No analog amplitude: the channel does not generate a continuous $y$ with noise; it directly flips bits.
        \item No natural reliability distribution: any LLR/quantization derived afterwards is \textbf{not} AWGN-consistent unless explicitly modeled.
      \end{itemize}
    \end{block}
  \end{column}
  \begin{column}{0.40\textwidth}
    \centering
    \begin{tikzpicture}[x=1cm,y=1cm]
      \node[font=\small] at (0,1.35) {Codeword bits ($N$)};
      \draw[rounded corners,fill=gray!10] (-3.1,0.55) rectangle (3.1,1.05);
      \foreach \x/\lbl in {-2.9/0,-2.3/1,-1.7/2,-0.5/\cdots,1.7/N{-}2,2.3/N{-}1} {
        \node[font=\tiny] at (\x,1.12) {\lbl};
      }
      % bit cells
      \foreach \i in {-3,-2.4,-1.8,-1.2,-0.6,0,0.6,1.2,1.8,2.4} {
        \draw[white,line width=0.5pt] (\i,0.55) rectangle (\i+0.55,1.05);
      }
      % mark K error positions
      \foreach \i in {-2.7,-1.5,0.3,2.1} {
        \draw[red,very thick] (\i,0.62) -- (\i+0.45,0.98);
        \draw[red,very thick] (\i+0.45,0.62) -- (\i,0.98);
      }
      \node[font=\small] at (0,0.15) {$K$ random positions flipped};
      \draw[->,thick] (0,-0.1) -- (0,-0.55);
      \node[font=\small] at (0,-0.95) {Fixed $FBC=K$};
    \end{tikzpicture}
  \end{column}
\end{columns}
\end{frame}

\begin{frame}{AWGN: BPSK with Additive Gaussian Noise}
\vspace{-0.2cm}
\begin{columns}
  \begin{column}{0.56\textwidth}
    \begin{block}{Signal model (conceptual)}
      \begin{itemize}
        \item BPSK mapping: bit $b\in\{0,1\}$ maps to $x=-(2b-1)$, so $b=0\to x=+1$ and $b=1\to x=-1$.
        \item Channel: $y=x+\sigma n$ with $n\sim\mathcal{N}(0,1)$ i.i.d.
        \item Hard decision: $\hat b=\mathbf{1}\{y<0\}$; the raw error count $FBC$ is a \textbf{random variable}.
        \item Soft decision: quantize $y$ into bins (by Vref) and look up LLR; reliability comes from distance to decision boundaries.
      \end{itemize}
    \end{block}
    \begin{block}{Why $FBC$ is not fixed}
      \begin{itemize}
        \item Even with fixed $\sigma$, each frame draws a different noise vector, so $FBC$ fluctuates around its mean.
        \item Typical use: estimate FER by Monte Carlo averaging over these naturally occurring fluctuations.
      \end{itemize}
    \end{block}
  \end{column}
  \begin{column}{0.40\textwidth}
    \centering
    \begin{tikzpicture}[x=1cm,y=1cm]
      \draw[->] (-3.2,0) -- (3.2,0) node[right] {$y$};
      \draw[->] (0,0) -- (0,2.1) node[above] {PDF};
      \draw[dashed] (0,0) -- (0,1.95) node[above] {$0$};
      % two Gaussians (stylized)
      \draw[thick,blue,domain=-3:3,samples=200] plot (\x,{1.8*exp(-(\x-1)^2/0.55)});
      \draw[thick,red,domain=-3:3,samples=200] plot (\x,{1.8*exp(-(\x+1)^2/0.55)});
      \node[blue,font=\small] at (1.6,1.6) {$b=0$};
      \node[red,font=\small] at (-1.6,1.6) {$b=1$};
      % show error regions
      \fill[blue,opacity=0.15] (-3,0) -- (-3,0.02) .. controls (-2.0,0.35) and (-1.0,1.2) .. (0,0) -- cycle;
      \fill[red,opacity=0.15] (0,0) -- (0,0) .. controls (1.0,1.2) and (2.0,0.35) .. (3,0.02) -- (3,0) -- cycle;
      \node[font=\scriptsize] at (-1.6,0.6) {errors for $b=0$};
      \node[font=\scriptsize] at (1.6,0.6) {errors for $b=1$};
      \node[font=\scriptsize] at (0,-0.45) {$FBC$ varies per frame};
    \end{tikzpicture}
  \end{column}
\end{columns}
\end{frame}

\begin{frame}{TAWGN (TRUNC\_AWGN): AWGN Conditioned on Fixed $FBC$}
\vspace{-0.2cm}
\centering
\begin{tikzpicture}[x=1cm,y=1cm]
  % distribution picture (centered)
  \draw[->] (-3.6,0) -- (3.6,0) node[right] {$y$};
  \draw[->] (0,0) -- (0,2.0) node[above] {PDF};
  \draw[dashed] (0,0) -- (0,1.85) node[above] {$0$};
  \draw[thick,blue,domain=-3.5:3.5,samples=220] plot (\x,{1.65*exp(-(\x-1)^2/0.60)});
  \node[blue,font=\small] at (1.7,1.55) {$\mathcal{N}(+1,\sigma^2)$};
  \fill[red,opacity=0.18] (-3.5,0) -- (-3.5,0.02) .. controls (-2.3,0.35) and (-1.0,1.1) .. (0,0) -- cycle;
  \fill[green!60!black,opacity=0.18] (0,0) -- (0,0) .. controls (1.0,1.1) and (2.3,0.35) .. (3.5,0.02) -- (3.5,0) -- cycle;
  \node[red,font=\scriptsize] at (-1.8,0.75) {error half-space};
  \node[green!60!black,font=\scriptsize] at (1.8,0.75) {correct half-space};

  % mask sketch (same slide)
  \draw[rounded corners,fill=gray!10] (-3.2,-0.75) rectangle (3.2,-0.25);
  \node[font=\scriptsize] at (-2.5,-0.50) {mask $S$ (size $k$)};
  \foreach \x in {-1.8,-1.0,-0.2,0.6,1.4,2.2} {
    \draw[white,line width=0.5pt] (\x,-0.75) rectangle (\x+0.6,-0.25);
  }
  \draw[red,very thick] (-1.0,-0.73) rectangle (-0.4,-0.27);
  \draw[red,very thick] (1.4,-0.73) rectangle (2.0,-0.27);
  \node[font=\scriptsize] at (0,-1.15) {exactly $k$ red slots $\Rightarrow$ fixed $FBC=k$};
\end{tikzpicture}

\vspace{-0.05cm}
\begin{itemize}
  \item Goal: generate an AWGN-like frame while enforcing a fixed hard-error count $k$ (per frame).
  \item Method: choose $S$ with $|S|=k$; for $i\in S$ draw $y_i$ from the \textbf{wrong} half-space, else draw from the \textbf{correct} half-space.
  \item Benefit: avoids slow reject-sampling (e.g., \texttt{continue} filters) but still preserves AWGN amplitudes under the condition $FBC=k$.
\end{itemize}
\end{frame}

\begin{frame}{TAWGN $\Leftrightarrow$ Rejection Sampling (1/3): Setup \& Method A}
\vspace{-0.15cm}
\begin{block}{Setup (assume all-zero codeword for clarity)}
\vspace{-0.1cm}
\begin{itemize}
  \item BPSK: transmit $x_i=+1$ and receive $y_i=x_i+n_i$, where $n_i\sim\mathcal{N}(0,\sigma^2)$ i.i.d.
  \item Hard-decision error at position $i$: $\hat b_i=1 \iff y_i<0 \iff n_i<-1$.
  \item Event $E_k$: the frame has exactly $k$ errors, i.e., $\left|\{i: n_i<-1\}\right|=k$.
\end{itemize}
\end{block}

\begin{block}{Method A: rejection sampling (generate then filter)}
\vspace{-0.1cm}
Draw a full noise vector $\bm{n}\sim f(\bm{n})=\prod_{i=1}^{N}\phi_\sigma(n_i)$; accept it only if $\bm{n}\in E_k$:
\[
f_A(\bm{n}) \;=\; f(\bm{n}\mid E_k)
 \;=\; \frac{f(\bm{n})\,\mathbf{1}\{\bm{n}\in E_k\}}{P(E_k)}.
\]
\end{block}
\end{frame}

\begin{frame}{TAWGN $\Leftrightarrow$ Rejection Sampling (2/3): Method B (construction)}
\vspace{-0.15cm}
\begin{block}{Method B: construct samples directly under $E_k$}
\vspace{-0.1cm}
\begin{enumerate}
  \item Choose an error index set $S_{\text{err}}\subset\{1,\dots,N\}$ uniformly at random, with $|S_{\text{err}}|=k$.
  \item For $i\in S_{\text{err}}$: draw $n_i$ from a truncated Gaussian on the \textbf{error} half-space $(-\infty,-1)$.
  \item For $j\notin S_{\text{err}}$: draw $n_j$ from a truncated Gaussian on the \textbf{correct} half-space $[-1,\infty)$.
\end{enumerate}
\end{block}

\begin{block}{Joint PDF of Method B}
\vspace{-0.1cm}
Let $p=P(n<-1)$ be the (single-bit) raw error probability. Then for a given $S_{\text{err}}$,
\[
f(\bm{n}\mid S_{\text{err}})
=
\Big(\prod_{i\in S_{\text{err}}}\frac{\phi_\sigma(n_i)}{p}\Big)
\Big(\prod_{j\notin S_{\text{err}}}\frac{\phi_\sigma(n_j)}{1-p}\Big)
\mathbf{1}\{\bm{n}\in E_k\}.
\]
With uniform set selection $P(S_{\text{err}})=1/\binom{N}{k}$, the unconditional PDF is
\[
f_B(\bm{n})
=
\frac{f(\bm{n})\,\mathbf{1}\{\bm{n}\in E_k\}}
{\binom{N}{k}\,p^k(1-p)^{N-k}}.
\]
\end{block}
\end{frame}

\begin{frame}{TAWGN $\Leftrightarrow$ Rejection Sampling (3/3): Equivalence \& Physical Intuition}
\vspace{-0.15cm}
\begin{columns}
  \begin{column}{0.56\textwidth}
    \begin{block}{Why $f_B(\bm{n})=f_A(\bm{n})$}
    \vspace{-0.15cm}
    \[
    P(E_k)=\binom{N}{k}p^k(1-p)^{N-k}\quad\Rightarrow\quad
    f_B(\bm{n})=\frac{f(\bm{n})\,\mathbf{1}\{\bm{n}\in E_k\}}{P(E_k)}=f(\bm{n}\mid E_k).
    \]
    \vspace{-0.1cm}
    \begin{itemize}
      \item Method A = ``sample then filter'' $\Rightarrow$ conditional distribution $f(\bm{n}\mid E_k)$.
      \item Method B = ``sample directly from $f(\bm{n}\mid E_k)$'' via truncated Gaussians.
      \item Therefore TAWGN (construction) is \textbf{exactly} rejection sampling, but without wasted trials.
    \end{itemize}
    \end{block}
  \end{column}
  \begin{column}{0.40\textwidth}
    \begin{block}{Physical/causal view: why it is valid}
    \vspace{-0.1cm}
    \begin{itemize}
      \item AWGN is \textbf{memoryless}: each $n_i$ only has a value, not a ``generation history''.
      \item Observing ``this bit is wrong'' simply means $n_i$ lies in the error half-space; its posterior is a truncated Gaussian.
      \item The decoder only sees $y$ (or LLR bins). It cannot tell whether $y$ came from filtering or direct construction.
      \item Method A acceptance rate is $P(E_k)$, so expected trials $\approx 1/P(E_k)$ (very slow at high SNR); Method B avoids this.
    \end{itemize}
    \end{block}
  \end{column}
\end{columns}
\end{frame}

\begin{frame}{High-Reliability Error (HRE): Definition and Why It Matters}
\vspace{-0.2cm}
\begin{columns}
  \begin{column}{0.56\textwidth}
    \begin{block}{Definition}
      \begin{itemize}
        \item HRE: a bit that is \textbf{actually wrong} but is reported with \textbf{very high confidence}.
        \item In quantized soft-decision, this corresponds to landing in an \textbf{extreme LLR bin} of the wrong side.
      \end{itemize}
    \end{block}
    \begin{block}{Two common categories}
      \begin{itemize}
        \item Fixed HRE: persistent defects or program/copy failures (bit flips are stable across reads).
        \item Probabilistic HRE: retention/disturb/read-offset tracking issues push some cells deep into the wrong high-confidence region.
      \end{itemize}
    \end{block}
    \begin{block}{Decoder impact (intuition)}
      \begin{itemize}
        \item LDPC relies on soft information; large-magnitude wrong LLRs can dominate message passing and amplify failures.
        \item Measuring both \textbf{FBC} and the \textbf{bin-level breakdown} helps explain when FER is driven by rare high-confidence wrong events.
      \end{itemize}
    \end{block}
  \end{column}
  \begin{column}{0.40\textwidth}
    \centering
    \fbox{\parbox[c][0.62\textheight][c]{0.95\linewidth}{
      \centering
      \textbf{(Figure placeholder)}\\[0.2cm]
      Paste your NAND $V_{th}$ distribution / read-threshold illustration here.
    }}
  \end{column}
\end{columns}
\end{frame}


{\aauwavesbg
\begin{frame}[plain,noframenumbering]
  \LARGE
  \finalpage{Thank you for listening!}
\end{frame}}
%%%%%%%%%%%%%%%%

\end{document}
